% =====================================================================
%  appendix_model_inputs.tex
%  Appendix A: record of every input value of the reactor model.
%  Values read from the simulation repository, branch main, commit
%  e4a3582 (8 September 2026). Formulas are given in the main text.
% =====================================================================

\chapter{Record of the model input data}
\label{app:model-inputs}
\label{app:labgene-data}

\section{Reading key}
\label{app:inputs-purpose}

This appendix records the value of every input of the reactor model,
its origin, and whether it was validated in this thesis. The formulas
that use these values are given in the main text. All values were read
from the simulation repository, branch \texttt{main}, commit
\texttt{e4a3582}, the state described in
Appendix~\ref{app:code-architecture}. The codes of the tables are given
in Table~\ref{tab:in-key}.

\begin{table}[htbp]
  \centering
  \small
  \caption{Codes of the input tables.}
  \label{tab:in-key}
  \begin{tabularx}{\textwidth}{@{}l>{\raggedright\arraybackslash}X@{}}
    \toprule
    Code & Meaning \\
    \midrule
    BRZ & Stated by the Brazilian Navy or the Brazilian government \\
    LIT & Stated in a peer-reviewed article, a laboratory report or a code manual \\
    REF & Taken from the NuScale-class reference benchmark adopted as design analog \\
    INF & Inferred from open information that does not state the value directly \\
    EST & Engineering estimate of this work \\
    NUM & Numerical setting with no physical counterpart \\
    DER & Derived from other entries \\
    \midrule
    V = Y & Confirmed by a verification calculation of this thesis \\
    V = P & Bounded by a sensitivity calculation of this thesis, not confirmed \\
    V = N & Adopted as given, not verified in this thesis \\
    \bottomrule
  \end{tabularx}
\end{table}

% =====================================================================
\section{Open data on the LABGENE prototype}
\label{app:inputs-labgene}

Table~\ref{tab:labgene-data} records the open information on the
LABGENE prototype used in this thesis. Four model inputs are anchored
to it: the reactor type, the thermal power, the enrichment class and the
vessel inner radius. Every other input is taken from the NuScale-class
reference benchmark (REF) or estimated in this work (EST).

\begin{table}[htbp]
  \centering
  \footnotesize
  \caption{Open-source data on the LABGENE prototype.}
  \label{tab:labgene-data}
  \begin{tabularx}{\textwidth}{@{}>{\raggedright\arraybackslash}p{3.2cm} >{\raggedright\arraybackslash}p{3.8cm} >{\raggedright\arraybackslash}X@{}}
    \toprule
    \textbf{Item} & \textbf{Value adopted} & \textbf{Source} \\
    \midrule
    Reactor type & PWR & \cite{ctmsp_labgene_2015,vasconcelos_fapesp_2025} \\
    Rated power & 48\,MW thermal & \cite{demartini_nuclep_2022}, 48\,MWth or 11\,MWe in \cite{galante_podernaval_2020}, 48 to 50\,MWth in \cite{dinizcosta_2017} \\
    Fuel & UO$_2$ pellets in rods grouped into fuel elements & \cite{marinha_inb_labgene_2022,pereira_rmb_2023} \\
    Enrichment class & LEU, below 20\,wt\% & \cite{brazil_npt_wp71_2022,brazil_npt_wp13_2026}, definition \cite{nrc_10cfr110_2} \\
    Enrichment value & Searched, up to 13.913\,wt\% in Campaigns 8 and 9 & 6 to 8\,wt\% \cite{dinizcosta_2017}, 18 to 19\,wt\% \cite{kassenova_2014}, up to 19\,wt\% \cite{vasconcelos_fapesp_2025} \\
    Refuelling interval & Objective in Campaigns 1 to 8, 1826\,EFPD constraint in Campaign 9 & Five years \cite{dinizcosta_2017} \\
    Vessel inner radius & 90\,cm & \cite[Fig.~13]{maruyama_2013} \\
    Vessel wall and height & 10\,cm and 470\,cm, not in the two-dimensional model & \cite[Fig.~13]{maruyama_2013} \\
    Active fuel height & 120\,cm, estimate & Proportions of \cite[Fig.~13]{maruyama_2013} \\
    Lattice, pitch, assembly count, absorbers, boron, reflector & Estimates of this thesis & No open source \\
    Plant status & Not critical as of May 2025, installations 66.8\,\% complete at the end of 2025 & \cite{vasconcelos_fapesp_2025,marinha_relatorio_gestao_2025} \\
    Safeguards status & Special procedures under Article~13 of the Quadripartite Agreement in negotiation & \cite{iaea_dg_statement_2022,eu_statement_bog_2025,brazil_npt_wp13_2026} \\
    \bottomrule
  \end{tabularx}
\end{table}

% =====================================================================

\section{Vessel, core envelope and reflector}
\label{app:inputs-vessel}

\small
\setlength{\tabcolsep}{3pt}
\begin{longtable}{@{}>{\raggedright\arraybackslash}p{3.25cm} >{\raggedright\arraybackslash}p{2.65cm} >{\raggedright\arraybackslash}p{1.15cm} >{\raggedright\arraybackslash}p{1.1cm} >{\raggedright\arraybackslash}p{3.85cm} >{\raggedright\arraybackslash}p{0.4cm} >{\raggedright\arraybackslash}p{0.95cm}@{}}
\caption{Vessel, core envelope and reflector.}
\label{tab:in-vessel} \\
\toprule
\textbf{Parameter} & \textbf{Value} & \textbf{Unit} & \textbf{Code} &
\textbf{Source} & \textbf{V} & \textbf{Note} \\
\midrule
\endfirsthead
\toprule
\textbf{Parameter} & \textbf{Value} & \textbf{Unit} & \textbf{Code} &
\textbf{Source} & \textbf{V} & \textbf{Note} \\
\midrule
\endhead
\bottomrule
\endlastfoot
Vessel inner radius & 90.0 & cm & INF & \cite[Fig.~13, p.~50]{maruyama_2013} & N & (a) \\
Vessel wall thickness & 10.0 & cm & INF & \cite[Fig.~13]{maruyama_2013} & N & (b) \\
Vessel height & 470 & cm & INF & \cite[Fig.~13]{maruyama_2013} & N & (b) \\
Core barrel thickness & 5.08 & cm & REF & \cite{fridman_nuscale_benchmark_2023} & P & \\
Downcomer allowance & 2.0 & cm & EST & This work & P & (c) \\
Reserved radial clearance & 7.08 & cm & DER & Barrel plus downcomer & Y & \\
Lattice-corner pad & 0.02 & cm & NUM & This work & Y & (d) \\
Core map & 32 assemblies, $6\times6$ minus corners & -- & INF & This work & P & (e) \\
Fuel envelope radius at 1.26\,cm & 77.231 & cm & DER & Equation~\eqref{eq:renv} & Y & \\
Largest reflector thickness at 1.26\,cm & 5.669 & cm & DER & Equation~\eqref{eq:ggeom} & Y & \\
Active fuel height & 120.0 & cm & INF & Proportions of \cite[Fig.~13]{maruyama_2013} & P & (f) \\
Radial reflector material & Homogenised SS304 and borated water & -- & EST & This work & P & \\
Reflector steel volume fraction & 0.90 & -- & EST & This work, 0.956 in \cite{ez_aldeen_2025} & P & (g) \\
\end{longtable}
\setlength{\tabcolsep}{6pt}
\normalsize

\noindent \textbf{Notes.}
\begin{enumerate}[label=(\alph*), nosep]
  \item Simplified structural drawing prepared for CTMSP, the only open geometric source on the prototype.
  \item Not in the two-dimensional model. Modelled only in the three-dimensional confirmation.
  \item Sensitivity to 3.7\,cm computed on the candidates only.
  \item Prevents an OpenMC surface from passing through a lattice corner. Defined once, in \texttt{core\_geometry.py}.
  \item Set by the heavy-metal mass consistent with the active height and the vessel.
  \item Enters the two-dimensional model only through the heavy-metal mass and the specific power.
  \item More moderating than the benchmark. Varied as an option in the three-dimensional confirmation.
\end{enumerate}

% =====================================================================
\section{Fuel assembly, fuel rod and guide tube}
\label{app:inputs-assembly}

\small
\setlength{\tabcolsep}{3pt}
\begin{longtable}{@{}>{\raggedright\arraybackslash}p{3.25cm} >{\raggedright\arraybackslash}p{2.65cm} >{\raggedright\arraybackslash}p{1.15cm} >{\raggedright\arraybackslash}p{1.1cm} >{\raggedright\arraybackslash}p{3.85cm} >{\raggedright\arraybackslash}p{0.4cm} >{\raggedright\arraybackslash}p{0.95cm}@{}}
\caption{Fuel assembly, fuel rod and guide tube.}
\label{tab:in-assembly} \\
\toprule
\textbf{Parameter} & \textbf{Value} & \textbf{Unit} & \textbf{Code} &
\textbf{Source} & \textbf{V} & \textbf{Note} \\
\midrule
\endfirsthead
\toprule
\textbf{Parameter} & \textbf{Value} & \textbf{Unit} & \textbf{Code} &
\textbf{Source} & \textbf{V} & \textbf{Note} \\
\midrule
\endhead
\bottomrule
\endlastfoot
Lattice & $17\times17$ & -- & REF & \cite{nuscale_htp2_2017,ez_aldeen_2025} & Y & \\
Fuel rods per assembly & 264 & -- & DER & $17^2$ minus 25 non-fuel positions & Y & \\
Guide tubes per assembly & 24 & -- & REF & Westinghouse $17\times17$ map & Y & \\
Instrumentation tubes per assembly & 1 & -- & REF & Westinghouse $17\times17$ map & Y & (a) \\
Fuel pellet outer radius & 0.406 & cm & REF & \cite{nuscale_htp2_2017,fridman_nuscale_benchmark_2023} & N & \\
Cladding inner radius & 0.414 & cm & REF & \cite{nuscale_htp2_2017,fridman_nuscale_benchmark_2023} & N & \\
Cladding outer radius & 0.475 & cm & REF & \cite{nuscale_htp2_2017,fridman_nuscale_benchmark_2023} & N & \\
Guide tube inner radius & 0.572 & cm & REF & \cite{nuscale_htp2_2017,fridman_nuscale_benchmark_2023} & N & \\
Guide tube outer radius & 0.612 & cm & REF & \cite{nuscale_htp2_2017,fridman_nuscale_benchmark_2023} & Y & (b) \\
Fuel rod pitch & 1.26 & cm & EST & Fixed from Campaign~8 & Y & (b) \\
Assembly pitch & 21.42 & cm & DER & $17\times$ rod pitch & Y & \\
Intra-assembly enrichment zones & Central $9\times9$ inner, remainder outer & -- & EST & This work & P & (c) \\
Gadolinia rod ladder & 12, 16, 20, 24, 32, 40 & rods & EST & After \cite{bae_gd_2023,epr_fsar_ch4,nuscale_htp2_2017} & P & \\
Spacer grid band height & 4.445 & cm & REF & \cite{fridman_nuscale_benchmark_2023} & P & (d) \\
Number of spacer grids & 4 & -- & EST & This work, 5 on 200\,cm in the benchmark & P & (d) \\
Spacer grid span & 41.4 & cm & EST & This work & P & (d) \\
First grid offset above fuel bottom & 3.555 & cm & REF & \cite{fridman_nuscale_benchmark_2023} & P & (d) \\
Grid strap inner box & 1.224 & cm & REF & \cite{fridman_nuscale_benchmark_2023} & Y & (b) \\
\end{longtable}
\setlength{\tabcolsep}{6pt}
\normalsize

\noindent \textbf{Notes.}
\begin{enumerate}[label=(\alph*), nosep]
  \item Modelled with the guide tube dimensions and material.
  \item Below a pitch of 1.224\,cm the lattice clips the guide-tube wall, up to 28\,\% at 1.150\,cm (Figure~\ref{fig:gt-clip}).
  \item Equal inner and outer enrichments from Campaign~8.
  \item Used only in the three-dimensional confirmation.
\end{enumerate}

% =====================================================================
\section{Material compositions and densities}
\label{app:inputs-materials}

\small
\setlength{\tabcolsep}{3pt}
\begin{longtable}{@{}>{\raggedright\arraybackslash}p{3.25cm} >{\raggedright\arraybackslash}p{2.65cm} >{\raggedright\arraybackslash}p{1.15cm} >{\raggedright\arraybackslash}p{1.1cm} >{\raggedright\arraybackslash}p{3.85cm} >{\raggedright\arraybackslash}p{0.4cm} >{\raggedright\arraybackslash}p{0.95cm}@{}}
\caption{Material compositions and densities.}
\label{tab:in-materials} \\
\toprule
\textbf{Parameter} & \textbf{Value} & \textbf{Unit} & \textbf{Code} &
\textbf{Source} & \textbf{V} & \textbf{Note} \\
\midrule
\endfirsthead
\toprule
\textbf{Parameter} & \textbf{Value} & \textbf{Unit} & \textbf{Code} &
\textbf{Source} & \textbf{V} & \textbf{Note} \\
\midrule
\endhead
\bottomrule
\endlastfoot
UO$_2$ density & 10.4 & g/cm$^3$ & REF & \cite{ezaldeen_zenodo_2025} & N & \\
(U,Gd)O$_2$ density & 10.2 & g/cm$^3$ & EST & This work & N & \\
$^{234}$U content & $0.007731\,w_{235}^{1.0837}$ & wt\% & LIT & \cite{bowman_suto_1996,scale_polaris_manual} & N & (a) \\
$^{236}$U content & 0 & wt\% & EST & Fresh uranium & N & \\
Gd$_2$O$_3$ density & 7.41 & g/cm$^3$ & LIT & Handbook value & N & \\
Gadolinium isotopic composition & Natural & -- & EST & This work & N & \\
Enrichment reduction of a gadolinia rod & 5\,\% relative per wt\% Gd$_2$O$_3$, minimum 0.2\,wt\% & -- & LIT & \cite{mildrum_segovia_2001}, minimum: this work & N & \\
Zircaloy-4 composition & Zr 98.2, Sn 1.45, Fe 0.20, Cr 0.10, O 0.05 & wt\% & REF & \cite{ezaldeen_zenodo_2025} & N & (b) \\
Zircaloy-4 density & 6.55 & g/cm$^3$ & REF & \cite{ezaldeen_zenodo_2025} & N & (b) \\
Helium gap density & 0.0015 & g/cm$^3$ & REF & \cite{ezaldeen_zenodo_2025} & N & \\
Coolant density & 0.72 & g/cm$^3$ & EST & This work & N & (c) \\
Coolant thermal scattering & \texttt{c\_H\_in\_H2O} & -- & LIT & ENDF/B-VII.1 & N & \\
Soluble boron & Natural boron, mass fraction of the water & -- & EST & This work & Y & \\
SS304 composition & Fe 69.5, Cr 19.0, Ni 9.5, Mn 2.0 & wt\% & REF & \cite{ezaldeen_zenodo_2025} & N & (b) \\
SS304 density & 7.90 & g/cm$^3$ & REF & \cite{ezaldeen_zenodo_2025} & N & (b) \\
Control rod number densities, at/(b$\cdot$cm) & From the benchmark scripts & -- & REF & \cite{ezaldeen_zenodo_2025} & N & \\
Cross-section library & ENDF/B-VII.1 & -- & LIT & \cite{chadwick_endf71_2011} & N & \\
Depletion chain & PWR chain for ENDF/B-VII.1 & -- & LIT & OpenMC, \nolinkurl{chain_endfb71_pwr.xml} & N & \\
\end{longtable}
\setlength{\tabcolsep}{6pt}
\normalsize

\noindent \textbf{Notes.}
\begin{enumerate}[label=(\alph*), nosep]
  \item Documented below 10\,wt\%. Applied up to the 16\,wt\% peripheral cap. An overestimate of $^{234}$U lowers $k$, the conservative direction.
  \item Not checked against a compendium of material compositions.
  \item IAPWS-IF97 gives 0.712\,g/cm$^3$ at 15.5\,MPa and 580\,K~\cite{wagner_iapws_if97_2000}, 1.1\,\% below the value used, a bias common to every design.
\end{enumerate}

% =====================================================================
\section{Operating conditions and power normalisation}
\label{app:inputs-operating}

\small
\setlength{\tabcolsep}{3pt}
\begin{longtable}{@{}>{\raggedright\arraybackslash}p{3.25cm} >{\raggedright\arraybackslash}p{2.65cm} >{\raggedright\arraybackslash}p{1.15cm} >{\raggedright\arraybackslash}p{1.1cm} >{\raggedright\arraybackslash}p{3.85cm} >{\raggedright\arraybackslash}p{0.4cm} >{\raggedright\arraybackslash}p{0.95cm}@{}}
\caption{Operating conditions and power normalisation.}
\label{tab:in-operating} \\
\toprule
\textbf{Parameter} & \textbf{Value} & \textbf{Unit} & \textbf{Code} &
\textbf{Source} & \textbf{V} & \textbf{Note} \\
\midrule
\endfirsthead
\toprule
\textbf{Parameter} & \textbf{Value} & \textbf{Unit} & \textbf{Code} &
\textbf{Source} & \textbf{V} & \textbf{Note} \\
\midrule
\endhead
\bottomrule
\endlastfoot
Core thermal power & 48.0 & MW & BRZ & \cite{demartini_nuclep_2022}, 48 to 50\,MWth in \cite{dinizcosta_2017} & N & \\
Number of assemblies & 32 & -- & INF & This work & P & \\
Fresh heavy-metal mass & 4808 & kg & DER & Pellet volume, density, U fraction, rod and assembly counts & Y & (a) \\
Uranium mass fraction of UO$_2$ & 0.8815 & -- & DER & Stoichiometry & Y & \\
Core specific power & 9.983 & W/\hspace{0pt}gHM & DER & Thermal power over heavy-metal mass & Y & \\
Fuel temperature & 900 & K & REF & \cite{ez_aldeen_2025} & N & \\
Cladding and structure temperature & 600 & K & REF & \cite{ez_aldeen_2025} & N & \\
Moderator temperature & 580 & K & EST & This work & N & (b) \\
System pressure & 15.5 & MPa & LIT & \cite{todreas_kazimi_2012,epr_fsar_ch4}, value of \texttt{mtc\_scan.py} & N & (c) \\
Soluble boron concentration & 1000 & ppm & EST & This work & P & \\
\end{longtable}
\setlength{\tabcolsep}{6pt}
\normalsize

\noindent \textbf{Notes.}
\begin{enumerate}[label=(\alph*), nosep]
  \item The 4733\,kg of the source comments is the earlier target mass. The model builds 4808\,kg.
  \item 20\,K below the 600\,K of the benchmark. The bias is common to every design and does not reorder the front.
  \item Not an input of the neutronic model. It enters only through the coolant density.
\end{enumerate}

% =====================================================================
\section{Control rod assemblies}
\label{app:inputs-cra}

\small
\setlength{\tabcolsep}{3pt}
\begin{longtable}{@{}>{\raggedright\arraybackslash}p{3.25cm} >{\raggedright\arraybackslash}p{2.65cm} >{\raggedright\arraybackslash}p{1.15cm} >{\raggedright\arraybackslash}p{1.1cm} >{\raggedright\arraybackslash}p{3.85cm} >{\raggedright\arraybackslash}p{0.4cm} >{\raggedright\arraybackslash}p{0.95cm}@{}}
\caption{Control rod assemblies.}
\label{tab:in-cra} \\
\toprule
\textbf{Parameter} & \textbf{Value} & \textbf{Unit} & \textbf{Code} &
\textbf{Source} & \textbf{V} & \textbf{Note} \\
\midrule
\endfirsthead
\toprule
\textbf{Parameter} & \textbf{Value} & \textbf{Unit} & \textbf{Code} &
\textbf{Source} & \textbf{V} & \textbf{Note} \\
\midrule
\endhead
\bottomrule
\endlastfoot
Rodded assemblies & 32 of 32 & -- & EST & This work & P & \\
Regulating banks & Inner 16 assemblies, RE1 to RE4 & -- & EST & This work & Y & \\
Shutdown banks & Outer 16 assemblies & -- & EST & This work & P & \\
Absorber of the screen & B$_4$C & -- & EST & This work & Y & \\
Alternative absorber & AIC, 80 Ag, 15 In, 5 Cd & wt\% & REF & \cite{fridman_nuscale_benchmark_2023} & P & (a) \\
Absorber radius, two-dimensional model & 0.433 & cm & REF & Benchmark notebook & N & (b) \\
Gap outer radius & 0.437 & cm & REF & \cite{fridman_nuscale_benchmark_2023,ezaldeen_zenodo_2025} & N & \\
Rod cladding outer radius & 0.484 & cm & REF & \cite{fridman_nuscale_benchmark_2023,ezaldeen_zenodo_2025} & N & \\
Rod cladding material & SS304L & -- & REF & \cite{fridman_nuscale_benchmark_2023} & N & \\
Absorber radii, 3D model & 0.427 AIC, 0.423 B$_4$C & cm & REF & \cite{fridman_nuscale_benchmark_2023} & P & \\
\end{longtable}
\setlength{\tabcolsep}{6pt}
\normalsize

\noindent \textbf{Notes.}
\begin{enumerate}[label=(\alph*), nosep]
  \item Used only in the three-dimensional confirmation, below the B$_4$C in the rod stack.
  \item One radius for both absorbers, between the two benchmark radii. Varied as an option in the three-dimensional confirmation.
\end{enumerate}

% =====================================================================
\section{Depletion schedule and transport settings}
\label{app:inputs-numerics}

\small
\setlength{\tabcolsep}{3pt}
\begin{longtable}{@{}>{\raggedright\arraybackslash}p{3.25cm} >{\raggedright\arraybackslash}p{2.65cm} >{\raggedright\arraybackslash}p{1.15cm} >{\raggedright\arraybackslash}p{1.1cm} >{\raggedright\arraybackslash}p{3.85cm} >{\raggedright\arraybackslash}p{0.4cm} >{\raggedright\arraybackslash}p{0.95cm}@{}}
\caption{Depletion schedule and transport settings.}
\label{tab:in-numerics} \\
\toprule
\textbf{Parameter} & \textbf{Value} & \textbf{Unit} & \textbf{Code} &
\textbf{Source} & \textbf{V} & \textbf{Note} \\
\midrule
\endfirsthead
\toprule
\textbf{Parameter} & \textbf{Value} & \textbf{Unit} & \textbf{Code} &
\textbf{Source} & \textbf{V} & \textbf{Note} \\
\midrule
\endhead
\bottomrule
\endlastfoot
Beginning-of-life burnup steps & 0.5, 1.0, 2.0, 4.0, 6.0 & MWd/\hspace{0pt}kgHM & EST & This work & Y & (a) \\
Burnup step after the first block & 4.0 & MWd/\hspace{0pt}kgHM & EST & This work & Y & \\
Steps per depletion restart & 2 & -- & NUM & This work & Y & (b) \\
Depletion horizon & 100.0 & MWd/\hspace{0pt}kgHM & NUM & This work & Y & \\
Depletion transport & 4000 particles, 60 batches, 20 inactive & -- & NUM & This work & P & \\
Core solve at beginning of life & 100\,000 particles, 170 batches, 60 inactive & -- & NUM & This work & P & \\
Shannon entropy mesh & $8\times8\times1$ & -- & NUM & This work & Y & \\
Depletion operator & \texttt{Coupled\hspace{0pt}Operator}, direct rate tallies & -- & NUM & OpenMC default~\cite{romano_2021_depletion} & N & \\
Depletion integrator & \texttt{Predictor\hspace{0pt}Integrator}, constant extrapolation & -- & NUM & OpenMC default~\cite{romano_2021_depletion} & N & (c) \\
Reaction-rate normalisation & Fission Q values of the depletion chain & -- & NUM & OpenMC default~\cite{openmc_docs} & N & \\
Matrix exponential & CRAM, order 48 & -- & NUM & OpenMC default~\cite{openmc_docs} & N & \\
Depletion power & Fresh heavy-metal mass $\times$ specific power, constant through the cycle & -- & NUM & This work & Y & \\
Temperature treatment & Interpolation, 294 to 1500\,K & -- & NUM & OpenMC setting & N & \\
Random seed & CRC32 hash of the design vector & -- & NUM & This work & Y & \\
\end{longtable}
\setlength{\tabcolsep}{6pt}
\normalsize

\noindent \textbf{Notes.}
\begin{enumerate}[label=(\alph*), nosep]
  \item Resolves the xenon, samarium and gadolinia transients, 13.5\,MWd/kgHM in total.
  \item Restart defect and correction in Section~\ref{sec:restart_defect}.
  \item First-order scheme of Equation~\eqref{eq:dep-ce}, so the accuracy is
        governed by the step lengths above.
\end{enumerate}

% =====================================================================
\section{Screening limits, zoning and derived targets}
\label{app:inputs-limits}

\small
\setlength{\tabcolsep}{3pt}
\begin{longtable}{@{}>{\raggedright\arraybackslash}p{3.25cm} >{\raggedright\arraybackslash}p{2.65cm} >{\raggedright\arraybackslash}p{1.15cm} >{\raggedright\arraybackslash}p{1.1cm} >{\raggedright\arraybackslash}p{3.85cm} >{\raggedright\arraybackslash}p{0.4cm} >{\raggedright\arraybackslash}p{0.95cm}@{}}
\caption{Screening limits, zoning and derived targets.}
\label{tab:in-limits} \\
\toprule
\textbf{Parameter} & \textbf{Value} & \textbf{Unit} & \textbf{Code} &
\textbf{Source} & \textbf{V} & \textbf{Note} \\
\midrule
\endfirsthead
\toprule
\textbf{Parameter} & \textbf{Value} & \textbf{Unit} & \textbf{Code} &
\textbf{Source} & \textbf{V} & \textbf{Note} \\
\midrule
\endhead
\bottomrule
\endlastfoot
Minimum $k_\mathrm{eff}^\mathrm{core}$ & 1.02 & -- & EST & This work & Y & \\
Maximum $k_\mathrm{eff}^\mathrm{core}$, Campaigns 8 and 9 & 1.166 & -- & DER & Section~\ref{sec:meth-ctrl} & Y & \\
Axial leakage factor & 1.0289 & -- & DER & Section~\ref{sec:meth-lax} & Y & \\
Radial peaking screen & 2.00, 1.65 in Campaign 9 & -- & EST, LIT & This work, \cite{robinson_gd10} & Y & \\
Controllability margin & 1000 & pcm & EST & This work & Y & \\
LEU screening cap & 19.75 & wt\% & LIT & \cite{jason_2016,onr_leu_2016} & N & \\
LEU definition & 20 & wt\% & LIT & \cite{nrc_10cfr110_2} & N & \\
Peripheral enrichment cap, Campaigns 8 and 9 & 16 & wt\% & EST & This work & Y & \\
Upper bound of the enrichment variable & 13.913 & wt\% & DER & $16/1.150$ & Y & \\
Lower bound of the enrichment variable & 3.0 in Campaign 8, 2.0 in Campaign 9 & wt\% & EST & This work & Y & \\
Central ring multiplier & 0.720 & -- & EST & Section~\ref{sec:meth-zoning} & Y & \\
Middle ring multiplier & 0.8933 & -- & DER & Fissile balance & Y & \\
Peripheral ring multiplier & 1.150 & -- & EST & Section~\ref{sec:meth-zoning} & Y & \\
Assemblies per ring & 4, 12, 16 & -- & DER & Figure~\ref{fig:zoning-rings} & Y & \\
Discharge-burnup screen & 55 & MWd/\hspace{0pt}kgHM & LIT & \cite{song_sanchez_2026} & N & \\
\end{longtable}
\setlength{\tabcolsep}{6pt}
\normalsize

% =====================================================================
\section{Deviations from the reference benchmark}
\label{app:inputs-departures}

\begin{enumerate}[nosep]
  \item Active height 120\,cm, against 200\,cm.
  \item Reflector steel volume fraction 0.90, against 0.956.
  \item One absorber radius of 0.433\,cm, against 0.423\,cm for B$_4$C and 0.427\,cm for AIC.
  \item Coolant density 0.72\,g/cm$^3$ at 580\,K, against 0.752\,g/cm$^3$ at 600\,K.
  \item Four spacer grids over 120\,cm, one in the plenum, against five over 200\,cm.
\end{enumerate}
